\section{``Love, and Do What You Will'': Augustine's Rule in Full}

\subsection{The exact sentence and the text it interprets}

The famous sentence is five Latin words: \latin{Dilige, et quod vis fac}---``Love, and do what you will'' (Augustine, \emph{In epistulam Ioannis ad Parthos tractatus} VII.8; PL 35, col.~2033). It is not ``do whatever you happen to want.'' It is not the longer hybrid ``do what thou wilt under love, against such there is no law.'' The latter combines an Augustinian idea with Galatians 5:23 and, in its wording, invites comparison with a later occult slogan. Precision restores the doctrine.

Augustine is preaching on 1 John 4. The love in question begins not with the self's preference but with God's initiative: God loved us and sent his Son; therefore we love one another. It is \latin{dilectio} or charity---the Spirit-given participation by which one adheres to God and wills the neighbor's true good. Its pattern is the self-offering of Christ, its object includes the concrete brother, and its fruit reveals whether one abides in God.

The imperative order is everything:

\begin{center}
\begin{tikzpicture}[node distance=5mm]
\node[guide emphasis, text width=.18\linewidth] (gift) {God first loves\\and gives the Spirit};
\node[guide node, right=of gift, text width=.18\linewidth] (root) {charity becomes\\the inward root};
\node[guide node, right=of root, text width=.18\linewidth] (prudence) {reason judges\\the fitting act};
\node[guide node, right=of prudence, text width=.18\linewidth] (fruit) {the will acts\\for true good};
\draw[guide arrow] (gift) -- (root);
\draw[guide arrow] (root) -- (prudence);
\draw[guide arrow] (prudence) -- (fruit);
\end{tikzpicture}
\end{center}

``Love'' is not a decorative condition added after sovereign choosing. It names the new root from which willing proceeds. ``Do what you will'' describes the breadth and spontaneity of a will healed at its source. The sentence is daring because charity is demanding, not because morality is thin.

\subsection{One surface, different acts}

Augustine prepares the maxim through moral contrasts. Father, Son, and Judas can each be said to hand Christ over; the outward phrase is similar, the intentions and formal acts utterly different (VII.7). A father's corrective severity and a kidnapper's caress can again reverse surface appearances (VII.8). One must therefore attend to the root, mind, and will, not judge moral quality from softness alone.

The lesson is not that intention magically changes every object. It is that a thin physical description does not yet identify a complete human act. ``Touching,'' ``striking,'' ``giving,'' ``withholding,'' and ``speaking'' can name materially similar motions whose objects, relationships, ends, and circumstances differ. Later Catholic moral analysis makes this explicit: object, intention, and circumstances must all be good; an evil object cannot be justified by a benevolent end (CCC 1750--1756).

Augustine's ancient example of paternal correction must also be read under justice and the dignity of the child. It cannot license rage, humiliation, disproportionate force, or abuse by calling them love. Augustine's own sequel says charity corrects ``without bitterness,'' loves what God made in the person, and opposes the error the sinner made (VII.11). Genuine charity seeks healing through a means prudence can judge proportionate; an abuser's self-description does not settle the matter.

\subsection{Charity is not sentiment and not sincerity alone}

Augustine explicitly warns that love is not tame listlessness. It can keep silence or cry out, correct or spare. The pairings matter because no single temperament monopolizes charity. A conflict-averse person cannot identify love with avoidance; an aggressive person cannot identify it with confrontation. Charity supplies the end; prudence, justice, courage, and temperance shape the act.

Nor is charity mere sincerity. One can sincerely desire a disordered good, sincerely mistake possession for care, or sincerely make another person serve one's own need. The theological virtue has an objective measure: God loved above all for his own sake and the neighbor loved in God. It fulfills the commandments because it adheres to the Good they express, not because a private feeling exempts the lover from them (Matt 22:34--40; Rom 13:8--10; ST II--II, q.~23, a.~8; q.~25, a.~1; q.~26, aa.~1--2).

This is why Benedict XVI could gloss Augustine's rule as the freedom of a will anchored in Christ and truth. One who loves Christ can do what he wills because, through love, his will is united with God's---not because autonomy has become absolute (First Vespers for Saints Peter and Paul, 28 June 2008). Freedom and responsibility remain indivisible.

\subsection{``Against such there is no law''}

The separate phrase belongs to Paul. After warning the Galatians not to use freedom as an opportunity for the flesh, he commands service through love and identifies the works of the flesh. The fruit of the Spirit is love, joy, peace, patience, kindness, goodness, fidelity, gentleness, and self-control; ``against such there is no law'' (Gal 5:13--23).

Paul is not abolishing moral norm. The Spirit produces what the law aims at and the flesh resists. Those who belong to Christ crucify disordered passions and walk by the Spirit (5:24--25). There is no law \emph{against} such fruit because the fruit is conformity to the divine good. The Spirit writes the law inwardly and gives its fulfillment the form of filial freedom.

Augustine and Paul can therefore be joined theologically, though not misquoted as one source:

\begin{studybox}{The synthesis}
Under charity, do what you will; for charity does not leave the will as it found it. It receives the will from the Creator, heals it by grace, binds it to truth, orders it through the virtues, and makes it fruitful in service. Against such fruit there is no law because divine love is the law's fulfillment, not its rival.
\end{studybox}

\subsection{A will under love is a will surrendered to Providence}

Charity unites the two halves of this article. Metaphysically, the will is a received secondary cause. Morally, its freedom is perfected by adherence to the supreme Good. Spiritually, it can entrust outcomes because its treasure is no longer identical with any finite outcome. The person whose root is charity can speak or be silent, act or wait, correct or spare, petition or accept---not arbitrarily, and without making self-protection the final norm.

The maxim is thus an examination of conscience before it is a permission: What do I love? What good am I actually choosing? Is the neighbor a person to be served or material for my project? Would I still call this love if truth frustrated my preference? Have I used Providence to escape duty, or control to escape trust? A will unable to submit these questions is not yet the will Augustine sets free.

At the deepest point, \latin{Dilige} is itself gift and command. God gives what he commands without acting in place of the lover. Grace plants and quickens the root; the person truly bears fruit. Here the mystery of Providence and freedom becomes life: ``work out your salvation'' because ``God is at work in you, both to will and to work'' (Phil 2:12--13).

\subsection{A later resemblance: what history does and does not establish}

Aleister Crowley's Thelemic formulas---``Do what thou wilt shall be the whole of the Law'' (\emph{Liber AL} I.40) and ``Love is the law, love under will'' (I.57)---are close enough in vocabulary to Augustine's maxim that a reader may reasonably ask about their relation. Resemblance, however, is not proof of direct descent. The checked evidence establishes a more limited history:

\begin{center}
\small
\begin{tabularx}{\linewidth}{@{}>{\raggedright\arraybackslash}p{.17\linewidth}>{\raggedright\arraybackslash}p{.25\linewidth}X@{}}
\toprule
Witness & Formula or reference & What it establishes \\
\midrule
Augustine, early fifth century & \latin{Dilige, et quod vis fac} & Charity received from God is the inward root from which the Christian wills and acts. \\
Rabelais, sixteenth century & The rule of the Abbey of Th\'el\`eme, conventionally rendered ``Do what you will'' & A major pre-Crowley occurrence of the will-formula within the wider literary genealogy. \\
Crowley, 1904 & \emph{Liber AL} I.40 and I.57 & The Thelemic formulas themselves; these passages do not name Augustine. \\
Crowley, 1926 & ``The Antecedents of Thelema'' & A retrospective acknowledgment: Crowley names Augustine, expressly says Augustine's context is ``by no means'' the meaning of \emph{Liber AL}, and calls Rabelais ``far more important.'' \\
\bottomrule
\end{tabularx}
\end{center}

Graham John Wheeler's history of the precept confirms that multiple genealogy and its changing uses. The evidence establishes a retrospective comparison, while Crowley's 1904 wording and the checked scholarship establish no direct borrowing from Augustine.

The theological difference can be stated without overclaiming the history. Crowley denied that ``will'' meant every passing appetite and interpreted it as a person's ``True Will.'' Yet his order still is not Augustine's Christian order. Augustine begins with the God who first loves, gives charity by grace, reveals the good, and thereby heals willing; Crowley's maxim places love ``under will'' within a Thelemic system that does not accept that Christian grammar of Creator, revelation, sin, grace, and beatitude. The two orders are therefore incompatible, and their word order expresses an inversion.

\begin{studybox}{The measured conclusion}
Historically: a documented retrospective comparison, not a proven direct textual corruption. Theologically: two incompatible orders of love and will. Augustine does not place love beneath autonomous willing; he places willing beneath charity, and charity beneath neither creature nor preference, because charity comes from God.
\end{studybox}
